\chapter{层流与湍流}

\intro{湍流也许是古典物理学中最后尚未解决的主要问题。Feynman称其为"经典物理学中最重要的未解问题"。}

\section{层流到湍流的转捩}

\subsection{Reynolds的经典实验}

1883年，Osborne Reynolds在曼彻斯特大学进行了著名的圆管流动转捩实验。他在水流入口处注入染色液，通过调节流量观察流态变化。Reynolds发现存在两个截然不同的流态：

\begin{itemize}
    \item \textbf{层流}：染色液呈细直线状，流体质点沿平行路径运动，相邻流层间仅有分子尺度的动量交换
    \item \textbf{湍流}：染色液迅速扩散至整个管截面，流体质点作不规则、三维、非定常的混沌运动，伴有强烈的横向混合
\end{itemize}

\begin{figure}[H]
\centering
\begin{tikzpicture}[scale=1.3]
    \definecolor{mypipe}{RGB}{70,70,70}
    \definecolor{myblue}{RGB}{30,120,200}
    \definecolor{myred}{RGB}{200,50,50}
    \definecolor{mylaminar}{RGB}{20,60,150}
    \definecolor{myturb}{RGB}{200,30,30}
    
    \fill[mypipe!20] (0,0.5) rectangle (6,3.5);
    \draw[mypipe, line width=2pt] (0,0.5) -- (6,0.5);
    \draw[mypipe, line width=2pt] (0,3.5) -- (6,3.5);
    
    \draw[myblue, line width=1pt] (0.5,2) -- (3,2);
    \node at (1.8,2.8) {\small 层流};
    
    \draw[->, line width=1pt] (6.2,2) -- (7.5,2);
    \node at (6.4,2.3) {$Re\uparrow$};
    
    \begin{scope}[shift={(0,-3.5)}]
        \fill[mypipe!20] (0,0.5) rectangle (6,3.5);
        \draw[mypipe, line width=2pt] (0,0.5) -- (6,0.5);
        \draw[mypipe, line width=2pt] (0,3.5) -- (6,3.5);
        
        \draw[myred!80, line width=0.8pt] (0.5,2) -- (1.2,2.1);
        \draw[myred!60, line width=0.6pt] (1.2,2.1) -- (1.6,1.9);
        \draw[myred!70, line width=0.6pt] (1.6,1.9) -- (2.1,2.2);
        \draw[myred!80, line width=0.7pt] (2.1,2.2) -- (2.7,1.7);
        \draw[myred!50, line width=0.5pt] (2.7,1.7) -- (3.0,2.4);
        \draw[myred!70, line width=0.6pt] (3.0,2.4) -- (3.5,1.8);
        \draw[myred!60, line width=0.6pt] (3.5,1.8) -- (4.0,2.1);
        \draw[myred!40, line width=0.4pt] (4.0,2.1) -- (4.3,2.6);
        \draw[myred!50, line width=0.5pt] (4.3,2.6) -- (4.7,1.6);
        \draw[myred!60, line width=0.5pt] (4.7,1.6) -- (5.2,2.3);
        \draw[myred!45, line width=0.4pt] (5.2,2.3) -- (5.5,1.7);
        \draw[myred!55, line width=0.5pt] (5.5,1.7) -- (5.8,2.2);
        \node at (3,2.9) {\small \hspace{0.5cm} 湍流};
    \end{scope}
\end{tikzpicture}
\caption{Reynolds实验示意：低Re时染色线保持清晰的直线（层流），超过临界Re后染色线迅速扩散（湍流）}
\label{fig:reynolds-exp}
\end{figure}

Reynolds进一步发现，层流向湍流转捩的临界Reynolds数 $Re_c \approx 2000$（圆管流动）。然而，若实验装置非常精细且来流无扰动，层流可持续到 $Re \approx 10^{5}$ 甚至更高。在实际工程中，$Re \approx 2300$ 通常取作转捩的工程判据。

\subsection{转捩机制}

层流到湍流的转捩并非单一的突变事件，而是一系列渐进的失稳过程。对于边界层和平板流动，转捩的经典路径包括：

\begin{enumerate}
    \item \textbf{感受性}：自由来流中的扰动（声波、涡扰动）激发边界层内的不稳定波
    \item \textbf{线性增长}：Tollmien-Schlichting波（T-S波）按线性稳定性理论指数增长
    \item \textbf{非线性饱和}：扰动振幅达到一定水平后非线性效应变得重要，出现二次失稳
    \item \textbf{三维化}：二维T-S波通过次谐波共振或K型/H型转捩发展为三维扰动
    \item \textbf{湍斑形成}：局部高剪切区出现Emmon湍斑（turbulent spots）
    \item \textbf{充分发展湍流}：湍斑合并，流动转变为充分发展的湍流
\end{enumerate}

\begin{remark}
上述"经典路径"适用于低来流湍流度条件。在来流湍流度较高时，扰动可以直接激发边界层的条带结构（streaks），绕过T-S波阶段直接转捩（bypass transition）。
\end{remark}

\section{流动稳定性理论}

\subsection{线性稳定性分析的基本框架}

流动稳定性理论（Hydrodynamic Stability Theory）研究基本流（层流）对小扰动的响应，是理解转捩的理论基础。将瞬时流动分解为基本流与扰动：

\begin{myformula}
    \mathbf{u}(\mathbf{x},t) = \mathbf{U}(\mathbf{x}) + \mathbf{u}'(\mathbf{x},t),\quad
    p(\mathbf{x},t) = P(\mathbf{x}) + p'(\mathbf{x},t)
\end{myformula}

代入NS方程，减去基本流满足的方程，并忽略二阶小量（线性化），得到线性化的扰动方程。对于平行剪切流 $\mathbf{U} = (U(y),0,0)$，引入流函数扰动 $\psi(x,y,t) = \phi(y)e^{i(\alpha x - \omega t)}$（法模态分析），得到：

\begin{myformula}
    \frac{1}{i\alpha Re}\left(\frac{d^{2}}{dy^{2}} - \alpha^{2}\right)^{2}\phi = (U-c)\left(\frac{d^{2}}{dy^{2}} - \alpha^{2}\right)\phi - \frac{d^{2}U}{dy^{2}}\phi
\end{myformula}

其中 $c = \omega/\alpha$ 为复数相速度。$c_i = \Im(c) > 0$ 对应时间增长（不稳定），$c_i < 0$ 对应衰减（稳定），$c_i = 0$ 对应中性稳定。

\subsection{Orr-Sommerfeld方程}

上述关于 $\phi(y)$ 的四阶常微分方程称为Orr-Sommerfeld方程（O-S方程），是线性流动稳定性理论的核心方程。方程中各项的物理意义：

\begin{itemize}
    \item $\frac{1}{Re}(D^{2}-\alpha^{2})^{2}\phi$：粘性扩散项（四阶导数，来自粘性项）
    \item $(U-c)(D^{2}-\alpha^{2})\phi$：惯性项（对流 + 非定常）
    \item $-U''\phi$：基本流涡度梯度的贡献
\end{itemize}

当 $Re \to \infty$ 时，忽略粘性项得到无粘稳定性方程（Rayleigh方程）：

\begin{myformula}
    (U-c)\left(\frac{d^{2}}{dy^{2}} - \alpha^{2}\right)\phi - \frac{d^{2}U}{dy^{2}}\phi = 0
\end{myformula}

\begin{myTheorem}{Rayleigh拐点定理}
平行剪切流 $\mathbf{U}=U(y)\mathbf{e}_x$ 无粘不稳定的必要条件是基本流速度剖面 $U(y)$ 存在拐点（即 $\frac{d^{2}U}{dy^{2}} = 0$ 在某点成立）。反之，若剖面无拐点，流动是无粘稳定的。
\end{myTheorem}

\begin{myTheorem}{Fjørtoft定理}
若速度剖面单调且无拐点，流动必是无粘稳定的。若存在拐点，则根据Fjørtoft的补充条件，$U''(U-U_s) < 0$ 在拐点 $y_s$ 附近必须成立（$U_s$ 为拐点处的速度），流动才是无粘不稳定的。
\end{myTheorem}

\begin{example}
Poiseuille流动 $U(y)=1-y^{2}$（$-1\le y\le 1$）在中心线 $y=0$ 处有拐点（$U''(0)=-2\neq 0$，但拐点要求 $U''=0$，故 $y=0$ 不是拐点——实际上 $U''=-2$ 处处不为零，故Poiseuille流无粘稳定），但粘性不稳定（Tollmien-Schlichting波）。Couette流动 $U(y)=y$ 处处无拐点（$U''=0$ 对任意 $y$ 成立，但按条件为零），既无粘稳定又粘性稳定，在所有Re下均为线性稳定。
\end{example}

\subsection{线性稳定性的数值求解}

O-S方程是四阶边值问题，需四个边界条件。对于槽道流动，通常取 $\phi(\pm1)=\phi'(\pm1)=0$（无滑移）。可采用Chebyshev谱方法高效求解。

\begin{mycode}{Python：Orr-Sommerfeld方程的特征值求解}
\begin{lstlisting}[language=Python]
import numpy as np
from scipy.linalg import eig
import matplotlib.pyplot as plt

def chebyshev_diff_matrices(N):
    """Build Chebyshev differentiation matrices of order 1 and 2."""
    points = np.cos(np.pi * np.arange(N+1) / N)
    c = np.ones(N+1); c[0] = 2; c[-1] = 2
    
    D1 = np.zeros((N+1, N+1))
    for i in range(N+1):
        for j in range(N+1):
            if i != j:
                D1[i,j] = (c[i]/c[j]) * (-1)**(i+j) / (points[i]-points[j])
    for i in range(N+1):
        D1[i,i] = -np.sum([D1[i,j] for j in range(N+1) if j != i])
    
    D2 = np.dot(D1, D1)
    D3 = np.dot(D2, D1)
    D4 = np.dot(D3, D1)
    return points, D1, D2, D3, D4

def solve_OS_plane_poiseuille(Re, alpha, N=80):
    """Solve Orr-Sommerfeld eq for plane Poiseuille flow."""
    y, D1, D2, D3, D4 = chebyshev_diff_matrices(N)
    y_inner = y[1:N]
    
    U = 1 - y_inner**2
    Upp = -2 * np.ones(N-1)
    
    I = np.eye(N-1)
    D_inner = D2[1:N, 1:N] - alpha**2 * I
    D2_inner = np.dot(D_inner, D_inner)
    
    A = D_inner / (1j * alpha * Re) + np.diag(U) @ D_inner - np.diag(Upp)
    B = D_inner
    
    eigenvalues, _ = eig(A, B)
    
    c = -1j * eigenvalues / alpha
    gamma_max = np.max(np.imag(c) * alpha)
    return y_inner, c, gamma_max

Re_vals = np.logspace(3, 4.5, 20)
alpha_vals = np.linspace(0.8, 1.2, 10)
neutral_curve = []

for Re_i in Re_vals:
    max_growth = -np.inf
    best_alpha = 0
    for alpha in alpha_vals:
        _, _, g = solve_OS_plane_poiseuille(Re_i, alpha, N=64)
        if g > max_growth:
            max_growth = g
            best_alpha = alpha
    if abs(max_growth) < 1e-4:
        neutral_curve.append((Re_i, best_alpha))

print("Neutral curve points computed.")
if neutral_curve:
    Rc, ac = zip(*neutral_curve)
    print(f"Critical Re approx: {min(Rc):.0f}")
\end{lstlisting}
\end{mycode}

\subsection{非线性稳定性与分岔}

线性理论只能预测扰动增长的初始阶段。当扰动振幅不再是小量时，非线性效应变得重要。对于二维流动，扰动演化可由Landau方程描述：

\begin{myformula}
    \frac{d|A|^{2}}{dt} = 2\sigma|A|^{2} - l|A|^{4}
\end{myformula}

其中 $A$ 为扰动振幅，$\sigma$ 为线性增长率，$l$ 为Landau常数（$l>0$ 对应超临界分岔，扰动饱和到有限振幅平衡；$l<0$ 对应亚临界分岔，需有限振幅扰动才能触发）。

\section{湍流的统计描述}

湍流的高度不规则性使得确定性描述极其困难。Reynolds（1895年）提出的统计方法是湍流理论的基础框架。

\subsection{Reynolds分解}

将瞬时流动变量分解为系综平均部分（时间平均部分）与脉动部分：

\begin{myformula}
    \mathbf{u}(\mathbf{x},t) = \bar{\mathbf{u}}(\mathbf{x}) + \mathbf{u}'(\mathbf{x},t),\quad
    p(\mathbf{x},t) = \bar{p}(\mathbf{x}) + p'(\mathbf{x},t)
\end{myformula}

对于定常湍流，时间平均定义为 $\bar{\mathbf{u}} = \frac{1}{T}\int_{0}^{T}\mathbf{u}\,dt$（$T$ 远大于湍流脉动的时间尺度）。平均运算满足以下性质（Reynolds条件）：

\begin{enumerate}
    \item $\overline{\bar{u}} = \bar{u}$（平均的幂等性）
    \item $\overline{u'} = 0$（脉动的平均为零）
    \item $\overline{\frac{\partial u}{\partial s}} = \frac{\partial\bar{u}}{\partial s}$（平均与微分可交换）
    \item $\overline{\bar{u}v'} = 0$（平均量与脉动量乘积的平均为零）
    \item $\overline{u'v'} \neq 0$（一般情形，关联量可非零）
\end{enumerate}

\begin{figure}[H]
\centering
\begin{tikzpicture}[scale=1.2]
    \definecolor{myinst}{RGB}{180,40,40}
    \definecolor{myavg}{RGB}{20,60,150}
    \definecolor{myfluct}{RGB}{60,60,60}
    
    \draw[->] (0,0) -- (8,0) node[right] {$t$};
    \draw[->] (0,-1.5) -- (0,2.5) node[above] {$u(t)$};
    
    \begin{scope}[
        /pgf/fpu=true,
        /pgf/fpu/output format=fixed,
    ]
    \draw[domain=0:7.5,samples=200,variable=\t,myinst,line width=0.8pt]
        plot ({\t}, {1.0 + 0.6*sin(25*\t r) + 0.3*sin(15*\t r)*cos(8*\t r)
        + 0.2*sin(40*\t r)});
    \end{scope}
    
    \draw[myavg, line width=2pt] (0,1.0) -- (7.5,1.0);
    \node[myavg] at (7.8,1.0) {$\bar{u}$};
    
    \draw[->] (5,2.2) -- (5.3,1.5) node[right] {\small 瞬时值 $u(t)$};
    
    \node[myfluct] at (4, -1.0) {$u(t)=\bar{u}+u'(t)$};
\end{tikzpicture}
\caption{Reynolds分解示意：瞬时速度分解为时间平均值与脉动值之和}
\label{fig:reynolds-decomp}
\end{figure}

\section{Reynolds平均Navier-Stokes方程}

\subsection{RANS方程的推导}

将Reynolds分解代入不可压缩NS方程并取系综平均，利用平均运算的性质，得到Reynolds平均Navier-Stokes方程（RANS方程）：

\begin{myformula}
    \frac{\partial \bar{u}_i}{\partial t} + \bar{u}_j\frac{\partial \bar{u}_i}{\partial x_j} = -\frac{1}{\rho}\frac{\partial \bar{p}}{\partial x_i} + \nu\frac{\partial^{2} \bar{u}_i}{\partial x_j\partial x_j} - \frac{\partial \overline{u_i'u_j'}}{\partial x_j}
\end{myformula}

与原始的NS方程相比，RANS方程右端多了一项 $-\frac{\partial\overline{u_i'u_j'}}{\partial x_j}$，该项来源于对流项中脉动量的非线性交互，称为Reynolds应力散度（或湍流应力散度）。

\begin{remark}
RANS方程是严格的，没有任何额外的近似。但RANS方程中出现了新的未知量——Reynolds应力 $\overline{u_i'u_j'}$（六个独立分量）。这导致著名的"封闭问题"：方程数少于未知量数。任何湍流模型都是在提供 $\overline{u_i'u_j'}$ 与平均量之间的关系来封闭RANS方程。
\end{remark}

\subsection{湍流的封闭问题}

RANS方程共有四个方程（连续性 + 三个动量方程），但未知量包括：
\begin{itemize}
    \item 平均速度 $\bar{\mathbf{u}}$（3个分量）
    \item 平均压力 $\bar{p}$（1个）
    \item Reynolds应力 $\overline{u_i'u_j'}$（6个独立分量）
\end{itemize}

共计10个未知量，而只有4个方程。这一问题被称为"湍流封闭问题"，是湍流建模的核心困难。各类湍流模型（代数模型、一方程模型、两方程模型、Reynolds应力模型、大涡模拟）本质上对应不同的封闭策略。

\section{Reynolds应力}

\subsection{物理意义}

$-\rho\overline{u_i'u_j'}$ 的物理意义是湍流脉动引起的动量输运。以剪切湍流为例，切向Reynolds应力 $-\rho\overline{u'v'}$ 表示 $x$ 方向动量沿 $y$ 方向的湍流输运率。

若平均速度梯度 $\frac{\partial\bar{u}}{\partial y} > 0$（如壁面附近），流体微团从高速区向低速区运动（$v' < 0$）时携带较高的 $x$ 方向动量（$u' > 0$），因此 $\overline{u'v'} < 0$，从而 $-\rho\overline{u'v'} > 0$，湍流应力的方向与粘性应力相同——促进动量输运。这一机理类似于分子粘性，但效率远高于分子过程。

\subsection{涡粘性假设}

Boussinesq（1877年）提出了最著名且至今仍被广泛使用的封闭假设——涡粘性模型：

\begin{myformula}
    -\rho\overline{u_i'u_j'} = \mu_t\left(\frac{\partial\bar{u}_i}{\partial x_j} + \frac{\partial\bar{u}_j}{\partial x_i}\right) - \frac{2}{3}\rho k\delta_{ij}
\end{myformula}

其中 $\mu_t$ 为涡粘性系数（湍流粘性），$k = \frac{1}{2}\overline{u_i'u_i'}$ 为湍动能。此假设将Reynolds应力与平均速度梯度线性关联，形式上与牛顿粘性定律一致，将封闭问题归结为确定 $\mu_t$（或等价地，湍流运动粘性 $\nu_t = \mu_t/\rho$）。

\subsection{混合长度模型}

Prandtl（1925年）提出了最早的涡粘性模型——混合长度模型。类比气体分子运动论中的平均自由程概念，假设湍流涡团在"混合长度" $\ell_m$ 内保持其动量特性：

\begin{myformula}
    \nu_t = \ell_m^{2}\left|\frac{\partial\bar{u}}{\partial y}\right|
\end{myformula}

混合长度 $\ell_m$ 依赖于距壁面的距离和流动构型。对于壁面湍流，$\ell_m = \kappa y$（近壁处，$\kappa \approx 0.41$ 为von Karman常数）。混合长度模型计算量小，物理直观，但需要预设 $\ell_m$ 的分布，对复杂流动（分离、回流）适用性有限。

\subsection{湍动能方程与两方程模型}

湍动能的精确输运方程可通过NS方程与RANS方程相减得到：

\begin{myformula}
    \frac{D k}{D t} = \mathcal{P} - \varepsilon + \frac{\partial}{\partial x_j}\left(\nu\frac{\partial k}{\partial x_j} - \frac{1}{2}\overline{u_i'u_i'u_j'} - \frac{1}{\rho}\overline{p'u_j'}\right)
\end{myformula}

各项物理意义：$\mathcal{P} = -\overline{u_i'u_j'}\frac{\partial\bar{u}_i}{\partial x_j}$ 为湍动能产生项（从平均流提取能量），$\varepsilon = \nu\overline{\frac{\partial u_i'}{\partial x_j}\frac{\partial u_i'}{\partial x_j}}$ 为湍动能耗散率。

$\varepsilon$ 本身也需要输运方程。最广泛使用的 $k$-$\varepsilon$ 模型（Launder \& Spalding, 1974）对 $\varepsilon$ 建模为：

\begin{myformula}
    \frac{D\varepsilon}{D t} = C_{\varepsilon 1}\frac{\varepsilon}{k}\mathcal{P} - C_{\varepsilon 2}\frac{\varepsilon^{2}}{k} + \frac{\partial}{\partial x_j}\left[\left(\nu + \frac{\nu_t}{\sigma_\varepsilon}\right)\frac{\partial\varepsilon}{\partial x_j}\right]
\end{myformula}

涡粘性由 $\nu_t = C_\mu k^{2}/\varepsilon$ 给出。标准常数：$C_\mu=0.09$，$C_{\varepsilon 1}=1.44$，$C_{\varepsilon 2}=1.92$，$\sigma_k=1.0$，$\sigma_\varepsilon=1.3$。

\section{湍流的能量级串——Kolmogorov理论}

\subsection{Richardson的级串概念}

Lewis Fry Richardson（1922年）以诗意的语言描述了湍流的本质：大涡产生小涡，小涡产生更小的涡，如此下去，直至粘度最终将动能耗散为热（"Big whorls have little whorls that feed on their velocity, and little whorls have lesser whorls, and so on to viscosity."）。

Richardson的级串图像描绘了能量从大尺度（含能尺度）通过惯性子区向小尺度（耗散尺度）逐级传递的过程。这一图像是Kolmogorov理论（K41理论）的物理基础。

\subsection{Kolmogorov假设}

Kolmogorov在1941年发表了三篇开创性论文，奠定了湍流统计理论的基础。其核心假设为局部各向同性假设：在足够高的Re数下，充分小的尺度上的湍流在统计上是各向同性的，与大尺度运动的各向异性细节无关。

定义如下特征尺度：
\begin{itemize}
    \item 含能尺度 $L$：能量注入的特征尺度，对应最大涡的尺寸
    \item Kolmogorov尺度（耗散尺度）$\eta$：粘性耗散主导的尺度
    \item Taylor微尺度 $\lambda$：介于含能尺度与耗散尺度之间的中间尺度
\end{itemize}

Kolmogorov尺度的量纲分析：
\begin{myformula}
    \eta = \left(\frac{\nu^{3}}{\varepsilon}\right)^{1/4},\quad
    u_\eta = (\nu\varepsilon)^{1/4},\quad
    \tau_\eta = \left(\frac{\nu}{\varepsilon}\right)^{1/2}
\end{myformula}

其中 $\varepsilon$ 为湍动能耗散率（单位质量的能量耗散率，量纲 $[L^{2}T^{-3}]$）。含能尺度与耗散尺度之比为 $\frac{L}{\eta} \sim Re^{3/4}$（基于 $L$ 尺度的 $Re$）。

\subsection{惯性子区与 $-5/3$ 律}

在含能尺度与耗散尺度之间存在"惯性子区"（$L \gg r \gg \eta$），其中粘性耗散可忽略，流动动力学由惯性（非线性对流）主导。Kolmogorov第二假设指出，惯性子区内的统计量仅依赖于 $\varepsilon$（能量通量）和尺度 $r$本身。

由量纲分析，结构函数（两点速度差的二阶矩）为：

\begin{myformula}
    \langle [u(x+r)-u(x)]^{2}\rangle = C_2 (\varepsilon r)^{2/3}
\end{myformula}

对应的能谱为著名的Kolmogorov $-5/3$ 律：

\begin{myformula}
    E(k) = C_K \varepsilon^{2/3} k^{-5/3}
\end{myformula}

其中 $C_K \approx 1.5$ 为Kolmogorov常数（由实验确定）。这一普适的 $-5/3$ 幂律已在大量实验和直接数值模拟（DNS）中得到验证，是湍流理论中最坚实的成果之一。

\begin{figure}[H]
\centering
\begin{tikzpicture}[scale=1.3]
    \definecolor{mycurve}{RGB}{20,60,150}
    \definecolor{my53}{RGB}{200,50,50}
    \definecolor{myinertial}{RGB}{200,180,220}
    
    \fill[myinertial!40] (2.5, 0) rectangle (7.5, 4);
    \node at (5, 3.8) {惯性子区};
    
    \draw[->, line width=1.5pt] (0,0) -- (12,0) node[below right] {$\log k$};
    \draw[->, line width=1.5pt] (0,0) -- (0,5) node[above] {$\log E(k)$};
    
    \draw[mycurve, line width=2pt] 
        (0.5,4.2) .. controls (2,3.5) and (3,2.5) .. (4,2.0)
        -- (8,0.4) .. controls (9.5,0.05) and (11,0) .. (11.5,0);
    
    \draw[my53, dashed, line width=1.5pt] (3.2,2.5) -- (8.2,1.0);
    \node[my53] at (9.2,1.8) {斜率 $-5/3$};
    
    \draw[decorate, decoration={brace, amplitude=8pt}] (0.3,4.2) -- (0.3,2.5)
        node[midway, left=10pt] {\small 含能区};
    \draw[decorate, decoration={brace, amplitude=8pt}] (9.5,0.5) -- (9.5,0)
        node[midway, right=10pt] {\small 耗散区};
    
    \draw[dashed] (2.2, 0) -- (2.2, 4);
    \draw[dashed] (8.2, 0) -- (8.2, 4);
    
    \node at (2.2, -0.4) {$k_L$};
    \node at (8.2, -0.4) {$k_\eta$};
\end{tikzpicture}
\caption{湍流能谱的Kolmogorov $-5/3$ 律示意：能谱在惯性子区内遵循 $k^{-5/3}$ 衰减规律}
\label{fig:kolmogorov-spectrum}
\end{figure}

\subsection{Kolmogorov常数与谱修正}

在过去的数十年里，大量实验与DNS研究验证了Kolmogorov $-5/3$ 律。Saddoughi \& Veeravalli（1994）在圆管、边界层和大气边界层等不同湍流中测量到 $C_K \approx 1.5\pm 0.1$ 的一致值。然而，精确的谱形状在含能—惯性子区过渡带和惯性子区—耗散区过渡带存在偏离，可以用对数修正或高阶展开描述。

值得注意的是，K41理论预测的精确 $-5/3$ 指数是基于量纲分析（假设惯性子区内唯一的决定参数是 $\varepsilon$）。Kolmogorov于1962年提出了修正理论（K62），引入了间歇性概念来校正高阶统计量的偏离。

\begin{mycode}{Python：模拟Kolmogorov能谱}
\begin{lstlisting}[language=Python]
import numpy as np
import matplotlib.pyplot as plt

def kolmogorov_spectrum(n_k=256, Re_lambda=200, C_K=1.5):
    """Generate a model Kolmogorov energy spectrum."""
    k = np.logspace(-1, 3, n_k)
    
    eta = 1.0
    L = eta * Re_lambda**(3/2)
    k_L = 2*np.pi / L
    k_eta = 2*np.pi / eta
    
    epsilon = 1.0
    E_k = C_K * epsilon**(2/3) * k**(-5/3)
    
    E_k *= np.exp(-1.5 * (k/k_eta)**(4/3))
    
    E_k *= (1 - np.exp(-(k/k_L)**4))
    
    return k, E_k, k_L, k_eta

k, E, k_L, k_eta = kolmogorov_spectrum()

fig, ax = plt.subplots(figsize=(8, 5))
ax.loglog(k, E, 'b-', linewidth=2, label='Model spectrum')
ax.loglog(k, 1.5 * k**(-5/3), 'r--', linewidth=1.5,
    label='$k^{-5/3}$')
ax.axvline(k_L, color='gray', linestyle=':', alpha=0.7)
ax.axvline(k_eta, color='gray', linestyle=':', alpha=0.7)
ax.text(k_L*3, 5e-6, 'Energy range', fontsize=10,
    rotation=90, va='bottom')
ax.text(k_eta/3, 5e-6, 'Dissipation range', fontsize=10,
    rotation=90, va='bottom')
ax.text(np.sqrt(k_L*k_eta), 2e-2, 'Inertial subrange',
    fontsize=10, ha='center')
ax.set_xlabel('$k$'); ax.set_ylabel('$E(k)$')
ax.legend(); ax.grid(True, alpha=0.3)
plt.tight_layout(); plt.show()
\end{lstlisting}
\end{mycode}

\section{湍流的间歇性}

\subsection{K41理论的局限性}

K41理论提供了一个优美且普适的湍流框架，但精确的实验测量揭示出与K41预测的系统偏差，尤其在更高阶统计量中。问题出在Kolmogorov假设的均匀性（self-similarity）：K41理论隐含假设了能量耗散率 $\varepsilon$ 在空间中均匀分布。然而，高Re数湍流中 $\varepsilon$ 是高度间歇的，耗散活动集中在空间的稀疏区域。

\subsection{间歇性的统计描述}

间歇性可通过高阶结构函数进行研究。第 $p$ 阶纵向速度结构函数定义为：

\begin{myformula}
    S_p(r) = \langle |u(x+r) - u(x)|^{p} \rangle
\end{myformula}

在惯性子区内，K41的量纲分析预测 $S_p(r) \sim (\varepsilon r)^{p/3} \sim r^{p/3}$，即标度指数 $\zeta_p = p/3$。但实验表明 $\zeta_p$ 与 $p/3$ 存在显著的偏离（异常标度），偏离程度随 $p$ 增大而增大。这被称为"湍流的异常标度率"。

\subsection{Kolmogorov-Obukhov修正（K62）}

Kolmogorov和Obukhov在1962年各自独立提出了修正框架。他们引入耗散率的空间平均 $\varepsilon_r$，假设 $\varepsilon_r$ 自身满足对数正态分布，且方差随尺度呈对数增长：

\begin{myformula}
    \langle (\ln\varepsilon_r)^{2} \rangle - \langle \ln\varepsilon_r \rangle^{2} = A + \mu\ln(L/r)
\end{myformula}

其中 $\mu$ 为间歇性指数（实验值 $\mu \approx 0.25$）。修正后的结构函数标度指数变为 $\zeta_p = p/3 - (\mu/18)p(p-3)$，引入了对 $p$ 的非线性依赖。

\subsection{多重分形模型与当前理解}

更现代的描述采用多重分形形式（Parisi \& Frisch, 1985）：湍流可视为具有多分形奇异性谱的随机场。不同的 $r$ 尺度上，速度增量的标度由Hölder指数 $h$ 表征：$\delta u(r) \sim r^{h}$，其中 $h$ 在不同位置取不同值，其概率密度服从多分形分布。多分形模型的标度指数为 $\zeta_p = \inf_h[ph + 3 - D(h)]$，其中 $D(h)$ 为 $h$ 对应的奇异性谱。

湍流间歇性仍是活跃的研究领域。高Re下的超级计算机DNS、实验室高精度流动测量和大气/海洋现场观测提供了丰富的数据，而这些数据的统计特性继续挑战着理论的完备性。

\begin{figure}[H]
\centering
\begin{tikzpicture}[scale=1.5]
    \definecolor{myK41}{RGB}{20,60,150}
    \definecolor{myK62}{RGB}{200,50,50}
    \definecolor{mymeas}{RGB}{50,150,50}
    
    \draw[->] (0,0) -- (6,0) node[right] {$p$};
    \draw[->] (0,0) -- (0,4.5) node[above] {$\zeta_p$};
    
    \draw[myK41, line width=2pt] (0,0) -- (5.8, 5.8/3);
    \node[myK41] at (5.2, 1.6) {K41: $\zeta_p=p/3$};
    
    \draw[myK62, line width=2pt] plot[domain=0:5.5,samples=100]
        ({\x}, {\x/3 - 0.25/18*\x*(\x-3)});
    \node[myK62] at (5.5, 2.7) {K62};
    
    \filldraw[mymeas] (1,0.36) circle (2pt);
    \filldraw[mymeas] (2,0.7) circle (2pt);
    \filldraw[mymeas] (3,1.0) circle (2pt);
    \filldraw[mymeas] (4,1.25) circle (2pt);
    \filldraw[mymeas] (5,1.5) circle (2pt);
    \node[mymeas] at (3.5, 3.8) {实验数据};
    
    \node at (5.8, -0.3) {};
\end{tikzpicture}
\caption{湍流结构函数的标度指数：K41线性预测与实验数据及K62修正的比较}
\label{fig:scaling-exponents}
\end{figure}

\section{湍流模型的工程实践}

尽管RANS湍流模型在数学上"不完美"（涡粘性假设缺乏严格的理论基础），它们在工程实践中取得了巨大成功。以下是工程中广泛使用的模型层级：

\subsection{零方程模型与一方程模型}

\textbf{零方程模型}（如Baldwin-Lomax模型）：完全用代数关系确定 $\nu_t$，无需求解输运方程。对附着湍流（边界层、槽道流）效果良好，但无法处理分离和回流。

\textbf{一方程模型}（如Spalart-Allmaras模型，SA模型）：求解单个关于湍流粘性的输运方程。SA模型在航空航天外流计算中非常成功，但对自由剪切流的预测能力有限。

\subsection{两方程模型}

$k$-$\varepsilon$ 模型及其变体（RNG $k$-$\varepsilon$，Realizable $k$-$\varepsilon$）是最广泛使用的湍流模型。$k$-$\omega$ 模型（Wilcox, 1988）用 $\omega \propto \varepsilon/k$（比耗散率）替代了 $\varepsilon$ 方程，在近壁处理上有优势。SST $k$-$\omega$ 模型（Menter, 1994）巧妙地将 $k$-$\omega$ 的近壁优点与 $k$-$\varepsilon$ 的远场优点结合，成为工业CFD中最受欢迎的RANS模型。

\subsection{Reynolds应力模型与高级方法}

二阶矩封闭模型（Reynolds Stress Model, RSM）直接求解Reynolds应力 $\overline{u_i'u_j'}$ 的六个输运方程，避免了涡粘性假设。RSM对旋转流、二次流和各向异性湍流的预测能力优于两方程模型，但计算成本高且数值稳定性差。

\textbf{大涡模拟（LES）}：显式解析大尺度的湍流运动，对小尺度（亚格子尺度）使用模型。LES的核心思想是：小尺度运动具有较多普适性和各向同性，比大尺度运动更容易建模。亚格子应力模型包括Smagorinsky模型（$\nu_{SGS} = (C_S\Delta)^{2}|\bar{S}|$）、动态模型（Germano, 1991）等。

\textbf{直接数值模拟（DNS）}：不对湍流作任何建模，直接离散求解NS方程。DNS是最精确但代价最高的方法。计算网格数需满足 $N^{3} \propto Re^{9/4}$，因此在可预见的未来只能用于低到中等Re数的湍流。

\begin{figure}[H]
\centering
\begin{tikzpicture}[scale=1.2]
    \definecolor{mydns}{RGB}{180,40,40}
    \definecolor{myles}{RGB}{40,120,40}
    \definecolor{myrans}{RGB}{40,40,180}
    \definecolor{mygrid}{RGB}{100,100,100}
    
    \draw[mygrid!20] (0,0) grid[xstep=0.5,ystep=0.5] (5,5);
    
    \foreach \x in {0,0.5,...,5} {
        \foreach \y in {0,0.5,...,5} {
            \draw[->, mydns, opacity=0.4] (\x,\y) -- ++(0.15,0.2);
        }
    }
    \node at (2.5, -0.5) {DNS: 所有尺度解析};
    
    \begin{scope}[shift={(7,0)}]
        \draw[mygrid!20] (0,0) grid[xstep=1.0,ystep=1.0] (5,5);
        \foreach \x in {0,1,...,5} {
            \foreach \y in {0,1,...,5} {
                \draw[->, myles, opacity=0.5] (\x,\y) -- ++(0.23,0.25);
            }
        }
        \node at (2.5, -0.5) {LES: 大尺度解析};
    \end{scope}
    
    \begin{scope}[shift={(14,0)}]
        \draw[mygrid!20] (0,0) grid[xstep=5,ystep=5] (5,5);
        \draw[->, myrans, line width=3pt, opacity=0.6] (2.5,2.5) -- (2.9,2.9);
        \node at (2.5, -0.5) {RANS: 全部建模};
    \end{scope}
\end{tikzpicture}
\caption{DNS、LES和RANS三种湍流模拟方法的比较示意}
\label{fig:turbulence-methods}
\end{figure}

\section{本章总结}

本章从Reynolds的经典转捩实验出发，系统阐述了流动从层流失去稳定性到发展为充分发展湍流的物理过程。线性稳定性理论（Orr-Sommerfeld方程）提供了转捩预测的理论框架。湍流的统计描述以Reynolds平均为基础，导出了RANS方程并引出了封闭问题。Kolmogorov的级串理论与 $-5/3$ 能谱律是湍流统计理论的基石，而间歇性修正揭示了真实湍流偏离K41预测的本质。在工程层面，各类湍流模型（从简单的混合长度模型到复杂的LES和DNS）形成了实践的层级，各自在精度与计算成本之间取得不同的平衡。
